\begingroup
\setlength{\parindent}{0pt}
\setlength{\parskip}{0.45\baselineskip}
\setlength{\abovedisplayskip}{7pt plus 1pt minus 1pt}
\setlength{\belowdisplayskip}{7pt plus 1pt minus 1pt}
\setlength{\jot}{5pt}
\clubpenalty=10000
\widowpenalty=10000
\section*{How the Proof Works}
\[
 \partial_tu+(u\cdot\nabla)u=-\nabla p+\nu\Delta u.
\]

The problem is to show that a smooth, incompressible three-dimensional flow
of finite energy stays smooth for all time when its viscosity is positive.
The difficulty comes from the fluid's ability to concentrate its motion.

When strain in the surrounding flow stretches a vortex, its material tube
lengthens and narrows. With circulation retained, that narrowing makes the
fluid spin faster. Viscosity spreads the rotation into the surrounding fluid,
but the surrounding strain can keep pulling the tube thinner. The tube can
even lose circulation while its rotation intensifies, if its cross-section
shrinks fast enough. Stretching and viscous spreading act together throughout
this process, and their relative rates decide whether the velocity differences
become sharper.

Continued concentration takes this competition into smaller and smaller
regions. Viscosity then has a shorter distance over which to smooth the flow,
but the faster motion can also leave it less time to do so. We need to know
whether viscosity gains enough ground as the tube narrows to prevent the
concentration from continuing all the way to a singularity.
The hope is simple: viscosity smooths the fluid faster than nonlinear
transport can sharpen it.

Rescaling lets us examine that hope while keeping viscosity fixed. For
\(\lambda>1\), the Navier--Stokes scaling is
\[
 u_\lambda(x,t)=\lambda u(\lambda x,\lambda^2t),
 \qquad p_\lambda(x,t)=\lambda^2p(\lambda x,\lambda^2t).
\]
Lengths contract, speeds increase, and the motion runs on a shorter clock.
Transport and diffusion preserve their relative strength. A smaller, faster
flow can still satisfy the equation with the original viscosity, so moving
to a finer scale gives diffusion no automatic advantage. At corresponding
times its total energy is smaller, too:
\(E[u_\lambda]=\lambda^{-1}E[u]\). Faster motion concentrated into less
space can cost less energy. A bound on total energy leaves that possibility
open.

Spatial Sobolev order tells us where a measurement begins to detect this
concentration. Increasing the order \(s\) gives more weight to finer spatial
variation, and under rescaling
\[
 \|u_\lambda(\cdot,t)\|_{\dot H^s}
 =\lambda^{s-1/2}\|u(\cdot,\lambda^2t)\|_{\dot H^s}.
\]
Below \(s=\tfrac12\), the norm decreases as the flow contracts. Above it,
the norm increases. At \(s=\tfrac12\), the measurement is unchanged.
This is the critical height: the flow can become arbitrarily small and fast
without changing its critical norm. Energy, at height zero, becomes smaller
throughout that comparison. Neither measurement by itself tells us how the
fluid will behave through a succession of such contractions.

The shorter clock is what makes a succession dangerous. If one contraction
could generate another, and each occurred \(\lambda^2\) times as quickly
as its predecessor, infinitely many stages would fit into a finite time:
\[
 \tau_j=\lambda^{-2j}\tau_0,
 \qquad \sum_{j=0}^{\infty}\tau_j
 =\frac{\tau_0}{1-\lambda^{-2}}<\infty.
\]
That is the hypothetical Zeno cascade. Rescaling describes flows with the
right spatial and temporal proportions for it. Sustaining the cascade would
require the evolving fluid to produce those stages, with each stage's
velocity, pressure and spatial structure generating what happens next.

For a singularity to form in finite time, velocity would have to blow up.
That requires acceleration to blow up, then jerk, and so on, and so on.
Here the temporal sequence is \(u,\partial_tu,\partial_t^2u,\ldots\),
evaluated at fixed spatial coordinates. A bounded next derivative keeps its
predecessor bounded on a finite interval. If the velocity becomes unbounded
in the spatial supremum norm as \(t\) approaches a finite \(T_*\), every
finite temporal order must participate in the escalation. 

But the equation makes each of these temporal responses depend on spatial
variation. The first response contains \(\nu\Delta u\); the next contains
\(\nu\Delta\partial_tu\), which reaches \(\nu^2\Delta^2u\) when the
equation is substituted. Continued differentiation keeps reaching further
into the spatial derivatives. Transport and pressure are differentiated
along with viscosity, so the connection includes how the whole velocity
field changes. The temporal escalation leads us into an increasingly
demanding spatial description of the fluid that would have to sustain it.

At the critical height, we can see this connection particularly clearly.
Write \(E_s=\tfrac12\|\Lambda^su\|_2^2\), with
\(\Lambda=(-\Delta)^{1/2}\), and set \(H=E_{1/2}\). The energy identity
at order \(s\) is
\[
 E_s'=-2\nu E_{s+1}+\mathcal P_s,
 \qquad
 \mathcal P_s=-\operatorname{Re}\langle\Lambda^su,
 \Lambda^s((u\cdot\nabla)u+\nabla p)\rangle.
\]
Viscosity relates change in time to energy one spatial height above. At the
critical height and the next temporal order, this gives
\[
 \begin{aligned}
 H'&=-2\nu E_{3/2}+\mathcal P_{1/2},\\
 H''&=4\nu^2E_{5/2}-2\nu\mathcal P_{3/2}
       +\partial_t\mathcal P_{1/2}.
 \end{aligned}
\]
A small \(H'\) can conceal a large viscous loss balanced by a large
nonlinear contribution. Once both terms are retained, their relation gives
\(E_{3/2}\) exactly. At the next order, the changing nonlinear contribution
must be retained as well; together the terms give \(E_{5/2}\). Every further
temporal order exposes another spatial height in this way. Positive viscosity
lets us recover each height because its coefficient \((-2\nu)^n\) is
nonzero.

The scale-invariant quantity has now led us to spatial measurements that
become progressively more sensitive to concentration. We reached them by
following its evolution, including the nonlinear changes that evolution
requires. Meanwhile, each temporal derivative of the velocity is itself a
spatial field: acceleration varies across the fluid, jerk varies across the
fluid, and their spatial derivatives remain related by the differentiated
equation. Temporal depth and spatial order belong together in a mixed
hierarchy.

This is where smoothness comes into view. Any finite Sobolev height gives
only a finite amount of spatial regularity. Yet for every finite derivative
order we want to control, a sufficiently high spatial Sobolev order supplies
it. Keeping all those orders together gives
\[
 L^2\cap\bigcap_{n\ge0}\dot H^{n+1/2}(\mathbb R^3)
 =\bigcap_{m\ge0}H^m(\mathbb R^3)
 =H^\infty(\mathbb R^3)\hookrightarrow C^\infty(\mathbb R^3).
\]
There is no final finite row at which all derivatives become controlled.
Smoothness belongs to the intersection. Each finite order has its own
requirement, and the velocity has to satisfy them together. The temporal
hierarchy is significant because the equation couples those spatial
requirements to the continuing evolution. If that coupled intersection
persists through the approach to \(T_*\), it gives both a smooth spatial
endpoint and the time continuity needed to reach it. The proposed singular
time becomes a time at which the fluid can continue.

If a singularity were to form, we would expect the opposite behavior.
The fluid would remain smooth before \(T_*\), while some fixed spatial
order lost control as \(T_*\) approached. Every earlier time could have
all its derivatives, yet the incoming evolution could fail to reach an
endpoint in the space those derivatives require. This is the distinction
that matters for the intersection: all orders must be retained through the
full approach to the proposed singular time. Differentiating a smooth
preterminal flow to arbitrary order establishes the structure; the proof
must establish its whole-interval realization.

We construct that realization from an initial velocity
\(u_0\in H^\infty_\sigma(\mathbb R^3)\). Its spatial derivatives are
available, and the equation fixes the first temporal response. Successive differentiation splits time derivatives across both
velocity factors in transport, with pressure changing alongside them:
\begin{flalign*}
 &\partial_tu=\nu\Delta u-(u\cdot\nabla)u-\nabla p,&&\\
 &\partial_t^2u=\nu\Delta\partial_tu-(\partial_tu\cdot\nabla)u
 -(u\cdot\nabla)\partial_tu-\nabla\partial_tp,&&\\
 &\partial_t^3u=\nu\Delta\partial_t^2u-(\partial_t^2u\cdot\nabla)u
 -2(\partial_tu\cdot\nabla)\partial_tu
 -(u\cdot\nabla)\partial_t^2u-\nabla\partial_t^2p.&&
\end{flalign*}
With \(V_n=\partial_t^nu\) and \(Q_n=\partial_t^np\), the pattern is
\[
 \begin{aligned}
 V_{n+1}&=\nu\Delta V_n
 -\sum_{a=0}^{n}\binom na(V_a\cdot\nabla)V_{n-a}-\nabla Q_n,\\
 -\Delta Q_n&=\partial_i\partial_j
 \sum_{a=0}^{n}\binom na V_{a,i}V_{n-a,j}.
 \end{aligned}
\]
The pressure row is determined by the velocity products, with pressure
normalized at spatial infinity. Evaluated at the initial datum, these
relations fix the temporal jet. Spatial differentiation fills out its rows.
As the fluid evolves, each row must continue satisfying these relations;
increasing the number of rows or their spatial order cannot change a
derivative already present.

The proof keeps finite pieces of this mixed structure. Choose temporal
depth \(N\) and spatial height \(M\). The corresponding rectangle contains
\(V_0,\ldots,V_N\), their adjacent temporal derivatives and their pressure
rows. The datum-generated rectangle theorem establishes, for every such
finite choice, the full-interval memberships
\[
 V_n\in L^2(0,T_*;H^{M+1}),\qquad
 V_{n+1}\in L^2(0,T_*;H^{M-1}),\qquad 0\le n\le N,
\]
with the pressure rows retained in \(L^1(0,T_*;H^M)\). The interval reaches
all the way to \(T_*\). Each rectangle contains derivatives of the original
datum-generated flow, and a larger rectangle restricts to the smaller one
where their coordinates overlap. This agreement lets the finite pieces
form one compatible intersection, \(\mathcal I[u_0]\). Every chosen
temporal and spatial order belongs to it, with no requirement for a common
bound across all orders.

We can now follow a temporal row to the endpoint with any chosen spatial
accuracy. To obtain an \(H^m\) limit for \(V_n\), take a rectangle high
enough that \(V_{n+1}\in L^2(0,T_*;H^m)\). The remaining change between
two late times satisfies
\[
 \|V_n(t)-V_n(s)\|_{H^m}
 \le (t-s)^{1/2}\|V_{n+1}\|_{L^2(s,t;H^m)},
 \qquad s<t<T_*.
\]
It tends to zero as both times approach \(T_*\). The row reaches an actual
limit in \(H^m\). A higher spatial choice gives a more regular limit whose
restriction agrees with the first. For the zeroth temporal row, these
compatible limits identify one velocity endpoint:
\[
 u(t)\longrightarrow u_*\ \text{in every }H^m\quad(t\uparrow T_*),
 \qquad u_*\in H^\infty_\sigma(\mathbb R^3)\hookrightarrow C^\infty.
\]
\needspace{7\baselineskip}
The higher temporal rows arrive with it. At whatever finite Sobolev height
a product or derivative needs, the recurrence and normalized pressure pass
to the limit and recover the endpoint jet from \(u_*\). Local existence
starts the fluid moving from this smooth datum; uniqueness joins that flow
to the incoming solution. The evolution continues beyond \(T_*\), so a
finite \(T_*\) cannot be maximal.

From the completed intersection, we derive the finite estimates and use
them. For any finite temporal order \(r\), spatial derivative order \(k\)
and time \(T\), choose \(m>k+\tfrac32\). Continuity in \(H^m\) and
Sobolev embedding give
\[
 \sup_{0\le t\le T}\|\partial_t^ru(t)\|_{W^{k,\infty}}
 \le C_{m,k}\sup_{0\le t\le T}\|\partial_t^ru(t)\|_{H^m}<\infty.
\]
In particular, the accumulated velocity gradient is finite. Return to the
stretching vortex: two distinct material points carried by the flow can
approach one another only as fast as that gradient allows. Their separation
obeys
\[
 |X(a,t)-X(b,t)|\ge |a-b|
 \exp\!\left(-\int_0^t\|\nabla u(s)\|_\infty\,ds\right)>0.
\]
On every finite interval, a positive material separation retains a positive
lower bound. The shrinking clocks in the scaling comparison allowed
infinitely many contractions to fit before a finite time. The actual flow's
accumulated strain is finite, so it cannot complete that collapse. The
estimates express, at each finite order where we need them, the regularity
already obtained from the intersection.
\endgroup
